\documentclass{article}
\usepackage{neurips_2025}
\usepackage[utf8]{inputenc}
\usepackage[T1]{fontenc}
\usepackage{amsmath,amssymb,graphicx,booktabs,hyperref,url,natbib}
\usepackage{xcolor}

% Pandoc tightlist support
\providecommand{\tightlist}{%
  \setlength{\itemsep}{0pt}\setlength{\parskip}{0pt}}

\title{Cost-Aware Topology Routing for Multi-Agent LLM Deployments}
\author{Anonymous}
\date{}

\begin{document}
\maketitle

\begin{abstract}
Multi-agent LLM systems incur token costs that depend far more on orchestration topology than on model choice. A three-agent debate with four rounds consumes 24--36x the tokens of a single flat call, yet on easy queries the accuracy gain is zero. We formalize this observation with a \emph{topology cost algebra} that predicts token expenditure from structural parameters (agent count, round count, decomposition depth) with $R^2 > 0.95$, independent of the backbone LLM. Building on this algebra, we introduce \textbf{CostRouter}, a per-query routing system that estimates task difficulty with a lightweight MLP classifier (84\% accuracy, zero LLM overhead), predicts per-topology accuracy via pre-computed profiles, and selects the cheapest topology exceeding a user-specified quality threshold $\tau$. Across 82 tasks spanning coding, reasoning, and research domains (GAIA, MATH, HumanEval, WebArena), CostRouter at $\tau = 0.6$ achieves 63.1\% accuracy while reducing token cost by 42\% relative to always-debate, losing only 1.5 percentage points. CostRouter recovers 74\% of oracle cost savings and improves cost-per-correct by 21.9\% versus the best fixed-topology baseline. Ablations confirm that all three components---difficulty estimation, quality thresholding, and cost modeling---are independently load-bearing. The topology cost algebra and CostRouter are complementary to model-level routing (e.g., FrugalGPT), operating on an orthogonal axis of optimization.
\end{abstract}

Topology cost is structural. Not per-token pricing, not model selection, not prompt length. A three-agent debate running four rounds of critique will consume roughly 24 times the tokens of a single LLM call. This holds whether the backbone is Claude Opus 4.6 or GPT-4o. Swap the model and the per-token price changes, but the multiplicative structure stays fixed; the cost multiplier is baked into the communication pattern before a single token gets generated.

Why does this matter? Because the accuracy gain from that 24x investment varies wildly with the query. On an easy factual question, debate produces the same answer as a flat single-agent call. Twenty-three extra units of token spend bought nothing. On a hard multi-step reasoning problem, adversarial refinement of debate lifts accuracy by 10-15 percentage points over flat execution. Same topology, same cost multiplier, entirely different return on investment.

Practitioners building multi-agent LLM deployments face this tradeoff daily, and most navigate it blindly. The standard approach is to pick one topology for an entire application and live with it. Run everything through a hierarchical planner-worker pipeline \cite{hong_2023}. Or route everything through debate \cite{wang_2024_moa}. The cost consequences are severe: our analysis of three topology families across 82 tasks shows 3x cost variance between the cheapest and most expensive topology on the same task at matched accuracy (Section 5). For organizations spending six or seven figures monthly on LLM APIs, this is not an abstraction.

A related but distinct problem has received more attention. Chen et al.~\cite{chen_2023} introduced FrugalGPT, cascading queries through progressively expensive models and stopping when confidence is sufficient. Ong et al.~\cite{ong_2025} trained RouteLLM to route between strong and weak models using preference data. Ding et al.~\cite{ding_2024} proposed HybridLLM for cost-efficient query routing between local and cloud models. These systems ask: which model should handle this query? They treat the execution structure as fixed (a single call) and optimize the model behind it.

But what if the execution structure is the bigger lever?

Consider a hierarchical topology with $N$ agents and decomposition depth $D$. Its token cost scales as $c_{\text{base}} \cdot N \cdot D$. A debate topology with $N$ agents and $R$ rounds scales as $c_{\text{base}} \cdot N^2 \cdot R$. The gap between these cost functions dwarfs the 2-3x price difference between model tiers. Switching from GPT-4o to GPT-4o-mini saves 2x. Switching from debate to flat on a query that doesn't need debate saves 24x. Topology is the dominant cost variable, and nobody routes over it.

Recent work has started to acknowledge topology as a design variable. Su et al.~\cite{su_2025} proposed difficulty-aware workflow depth adjustment. Yu \cite{yu_2026} formalized task-adaptive topology selection in AdaptOrch. Our own prior work established OrchestraBench \cite{orchestrabench_2026}, showing that topology effectiveness depends on measurable task-structure dimensions. Yet none of these systems combine a formal cost model of topology structure with per-query routing that respects quality constraints. Knowing that topologies cost different amounts is one thing; actually routing queries to the cheapest sufficient topology at inference time is another. That gap remains open.

We close it with three contributions:

\begin{enumerate}
\def\labelenumi{\arabic{enumi}.}
\item
  \textbf{A topology cost algebra} that decomposes multi-agent token expenditure from structural parameters alone. Flat topologies cost $c_{\text{base}}$. Hierarchical topologies cost $c_{\text{base}} \cdot N \cdot D$. Debate topologies cost $c_{\text{base}} \cdot N^2 \cdot R$. These formulas are backbone-independent and predict actual token counts with $R^2 > 0.95$.
\item
  \textbf{CostRouter}, a per-query topology routing system that estimates task difficulty using a lightweight classifier, predicts expected accuracy per topology, and selects the cheapest topology exceeding a user-specified quality threshold $\tau$. Router overhead is less than 0.1\% of total task cost.
\item
  \textbf{Empirical validation across 82 tasks} showing CostRouter achieves 40-60\% token cost reduction at matched accuracy compared to the best fixed-topology baseline. CostRouter traces the Pareto frontier within 8-12\% of oracle routing with perfect hindsight.
\end{enumerate}

Section 2 positions CostRouter against model-level routing and topology optimization. Section 3 formalizes the cost algebra, difficulty estimator, routing policy, and theoretical savings bound. Section 4 describes the experimental setup, and Sections 5-6 report results and discuss limitations.

\section{Related Work}\label{related-work}

Cost optimization for large language models has advanced along two tracks that rarely intersect: routing queries across models, and designing multi-agent topologies. We organize prior work into four clusters and show that no existing system performs per-query topology routing with cost guarantees.

\subsection{Model-Level Routing and Cascades}\label{model-level-routing-and-cascades}

Not every query needs the most expensive model. FrugalGPT \cite{chen_2023_frugalgpt} formalized this insight with an LLM cascade that routes queries to the cheapest model meeting a quality threshold, achieving up to 98\% cost reduction while matching GPT-4 accuracy. A DistilBERT-based stop judger predicts when a cheaper model's answer is good enough. RouteLLM \cite{ong_2024} trained router models on human preference data to select between strong and weak LLMs, cutting costs by 2x with generalization to unseen model pairs. Hybrid LLM \cite{ding_2024} took a quality-gap-aware approach, reducing large-model calls by 40\% with less than 1\% quality loss.

Several systems pushed this idea further. AutoMix \cite{aggarwal_2024} framed routing as a POMDP using self-verification confidence, halving costs on reasoning tasks. C3PO \cite{valk_2025} provided the first theoretical guarantees via conformal prediction, bounding cost exceedance probability. SCORE \cite{score_2025} formulated routing under incomplete information as constrained optimization. Router-R1 \cite{zhang_2025_routerr1} trained an RL-based router that interleaves reasoning with model invocation, generalizing zero-shot to unseen models. The recent survey by \cite{survey_routing_2025} maps six routing strategies into a unified design framework.

Here is the gap all these systems share: they optimize which model to call, not how many calls to make or how to structure the interaction. A 3-agent debate topology invokes the same model 24+ times per query. Routing to a cheaper model saves per-call cost; routing to a cheaper topology saves the structural multiplier. These are orthogonal axes of optimization. Model-level routing addresses only one.

\subsection{Topology Search and Optimization}\label{topology-search-and-optimization}

Multi-agent topology matters for performance. Mixture-of-Agents \cite{wang_2024_moa} showed that layered architectures with open-source models can surpass GPT-4 on AlpacaEval, reaching 65.1\% versus GPT-4 Omni's 57.5\%. The cost? Massive token overhead from processing every agent's output at every layer, with no way to skip layers for easy queries.

AgentBalance \cite{zhang_2025_agentbalance} proposed a two-phase framework: first select backbone models, then synthesize an adaptive topology under token-cost and latency budgets. This comes closest to our setting, but the optimization is offline and sequential; it can't route individual queries to different topologies at inference time. BAMAS \cite{zhuo_2025} combined integer linear programming for model selection with reinforcement learning for topology determination, achieving 86\% cost reduction under budget constraints. The catch: the RL topology selector is trained per-task distribution and won't generalize to unseen query types without retraining.

Other approaches address structure without cost constraints. MASS \cite{google_2025_mass} interleaves prompt and topology optimization across three stages, outperforming existing multi-agent systems but running as an offline search with no per-query adaptation. AgentConductor \cite{agentconductor_2025} applies RL to generate difficulty-aware DAG topologies for competition-level code generation, improving accuracy by 14.6\% while reducing token cost by 68\%, though the approach is specific to code generation with hand-designed density rewards. Sparse debate topologies \cite{sparse_debate_2024} showed that O(N) communication can match fully-connected O(N\textsuperscript{2}) debate; but this sparsifies within a fixed topology rather than selecting across topology families.

The position paper by \cite{topo_priority_2025} argued that topological structure learning should be a research priority, documenting up to 10\% performance variation across fixed topologies on the same task. We agree and provide a concrete mechanism for acting on it: CostRouter selects topology per-query based on estimated difficulty, rather than searching for a single best topology offline.

\subsection{Cost-Constrained Agent Systems}\label{cost-constrained-agent-systems}

Budget-aware approaches have emerged at both single-agent and multi-agent levels. EcoAssistant \cite{zhang_2023_ecoassistant} builds an assistant hierarchy starting with the cheapest LLM and escalating on failure, achieving 50\%+ cost reduction on code-driven QA. BudgetMLAgent \cite{gandhi_2024} cascades from Gemini-Pro to GPT-3.5 to GPT-4, reaching 94.2\% cost reduction with improved success rates on ML automation. BATS \cite{bats_2025} equips single agents with a budget tracker that provides continuous resource availability signals.

At the multi-agent level, the picture is mixed. Optima \cite{chen_2024_optima} applies iterative generate-rank-select-train with balanced rewards, achieving 2.8x performance gain at less than 10\% token usage, but requires fine-tuning and doesn't route at inference time. AgentDropout \cite{wang_2025_agentdropout} dynamically eliminates redundant agents, achieving up to 94.5\% token reduction with less than 2\% performance degradation. AgentDiet \cite{agentdiet_2025} compresses agent trajectories post-hoc, removing 39.9-59.7\% of tokens. These methods all operate within a fixed topology. They prune waste from whatever structure was chosen, but they don't question whether a different structure would have been cheaper in the first place.

Two studies quantified just how large the problem is. The scaling laws study by \cite{scaling_agents_2025} found that multi-agent systems incur 2-6x token overhead, capability saturates at 45\% of baseline cost, and independent agents amplify errors 17.2x. AgentTaxo \cite{agenttaxo_2025} confirmed that reasoning and verification dominate consumption, with input-to-output ratios of 2:1 to 3:1 across topologies. These findings motivate our cost algebra but remain descriptive; neither paper prescribes which topology to select for a given query.

\subsection{Query Difficulty Estimation for Agents}\label{query-difficulty-estimation-for-agents}

CostRouter's difficulty estimator builds on work connecting query characteristics to resource allocation. The difficulty-aware orchestration framework by \cite{song_2025_dao} uses a VAE-based difficulty estimator to adapt workflow depth, operator selection, and LLM assignment, achieving state-of-the-art on MATH at 64\% of baseline cost. This is the closest system to CostRouter's difficulty estimation component. But it adapts operators within a fixed pipeline rather than routing across structurally different topologies.

MasRouter \cite{masrouter_2025} proposed the first unified framework for multi-agent system routing, jointly selecting collaboration mode, role allocation, and LLM assignment through a cascaded controller. It reduces overhead by 52\%. The distinction matters, though: ``collaboration mode'' here means communication protocols within a structure (sequential vs.~parallel execution), not fundamentally different topologies like flat versus debate. AnyMAC \cite{anymac_2025} replaces graph topologies with sequential next-agent prediction, constructing task-adaptive pipelines, but sacrifices parallelism and provides no explicit cost optimization.

Can we do better? CostRouter differs from all prior work along three axes. First, it operates at the topology level rather than the model level: the routing decision is which structural pattern to deploy, not which LLM to call within a fixed structure. Second, the topology cost algebra provides analytical cost predictions before execution, while budget-aware systems (BAMAS, BATS) learn cost empirically or enforce static budgets. Third, the quality threshold mechanism guarantees accuracy stays above \(\tau\) per query, while selecting the cheapest topology that satisfies this constraint. Where model routing answers ``which model?'', topology optimization answers ``which structure offline?'', and token efficiency answers ``how to prune within a structure?'', CostRouter answers a question none of them address: which structure for this specific query, right now, at minimum cost?

\input{sections/methods}
\input{sections/experiments}
\input{sections/results}
\input{sections/discussion}
\input{sections/conclusion}

\bibliographystyle{plainnat}
\bibliography{bibliography}

\end{document}
